\documentclass[11pt]{article}
\usepackage[margin=1in]{geometry}
\usepackage{xcolor}
\usepackage{hyperref}
\definecolor{astateal}{HTML}{0A3235}
\definecolor{astagreen}{HTML}{0FCB8C}
\hypersetup{colorlinks=true,linkcolor=astateal,urlcolor=astateal,citecolor=astateal}
\title{A small paper in an Asta workspace}
\author{Workspace example}
\date{}

\begin{document}
\maketitle

\section{A paper paired with Overleaf}
This paper is paired with an Overleaf project, so co-authors can edit it there while the workspace reviews every change in a pull request. Overleaf sees only this \texttt{paper} directory, so the paper keeps its own \texttt{references.bib} here, which co-authors can edit too. Vaswani and colleagues proposed the Transformer architecture based on attention mechanisms~\cite{vaswani2017attention}. That entry was copied from the workspace's \texttt{references.bib}, which the companion Quarto page cites.

\section{A compact figure}
Figure~\ref{fig:workspace} shows where each writing format gets its citations.
\begin{figure}[htbp]
  \centering
  \fcolorbox{astateal}{astagreen!15}{\parbox{0.72\linewidth}{\centering\vspace{0.7em}Quarto page \quad $\longleftarrow$ \quad \textbf{references.bib}\\[0.5em]copied entries \quad $\longrightarrow$ \quad \textbf{paper/references.bib} \quad $\longleftrightarrow$ \quad Overleaf\vspace{0.7em}}}
  \caption{The workspace bibliography feeds the Quarto page directly and the paired paper through copied entries.}
  \label{fig:workspace}
\end{figure}

\bibliographystyle{plain}
\bibliography{references}
\end{document}
